% ===========================================================================
%  Chapitre 5 — Fonctions logarithmes
%  PE 2026, composante ANALYSE
% ===========================================================================
\bmtchapitre{Fonctions logarithmes}
  {Étudier et représenter la fonction logarithme népérien ; résoudre
   équations, inéquations et systèmes faisant intervenir le logarithme ;
   utiliser le logarithme décimal et le logarithme de base quelconque.}
  {Analyse \textbullet\ environ 10 heures}

Le chapitre précédent laisse une lacune, et elle est visible : parmi les
primitives usuelles, celle de $x \mapsto \dfrac{1}{x}$ manque. Toutes les
fonctions puissances $x^{n}$ ont pour primitive $\dfrac{x^{n+1}}{n+1}$, sauf
une, celle où $n = -1$, car la formule y produirait une division par zéro.

Cette exception n'est pas un accident de calcul. Elle signale qu'aucune des
fonctions connues jusqu'ici ne convient : il faut en construire une nouvelle.
C'est ce que fait ce chapitre. La fonction obtenue, le logarithme népérien,
transforme les produits en sommes — propriété qui, avant les calculatrices,
servait à multiplier de grands nombres, et qui sert aujourd'hui à résoudre des
équations où l'inconnue est en exposant.

% ---------------------------------------------------------------------------
\begin{revision}
Aucune ligne sur le logarithme : uniquement ce qui va servir à le construire.

\begin{enumerate}[label=\textbf{\arabic*.}, leftmargin=*, itemsep=0.35em]
  \item Qu'est-ce qu'une primitive d'une fonction $f$ sur un intervalle ?
  \item Énoncer le théorème des valeurs intermédiaires dans le cas d'une
        fonction strictement monotone.
  \item Rappeler la dérivée de $x \mapsto \dfrac{1}{x}$ et son ensemble de
        définition.
  \item Que signifie « $f$ réalise une bijection de $I$ sur $J$ » ?
  \item Donner la dérivée de la réciproque d'une bijection dérivable.
  \item Résoudre $\dfrac{2x-1}{x+3} > 0$.
\end{enumerate}
\end{revision}

\begin{automatismes}
Sans calculatrice, sans brouillon.

\begin{enumerate}[label=\textbf{\alph*)}, leftmargin=*, itemsep=0.25em]
  \item Dérivée de $x \mapsto \dfrac{1}{x}$
  \item Une primitive de $x \mapsto x^{-3}$
  \item $\limpinf \dfrac{1}{x}$
  \item $\displaystyle\lim_{x \to 0^{+}} \dfrac{1}{x}$
  \item Signe de $\dfrac{x-2}{x}$ sur $\Rps$
  \item Dérivée de $x \mapsto 3x^{2}-1$
  \item Résoudre $x^{2}-3x+2 > 0$
  \item Ensemble de définition de $x \mapsto \dfrac{1}{x-5}$
\end{enumerate}
\end{automatismes}

\begin{corrige}[automatismes]
a) $-\dfrac{1}{x^{2}}$ \quad b) $-\dfrac{1}{2x^{2}}$ \quad c) $0$ \quad
d) $+\infty$ \quad e) négatif sur $\intervalleo{0}{2}$, positif au-delà \quad
f) $6x$ \quad g) $x<1$ ou $x>2$ \quad h) $\R \setminus \{5\}$.
\end{corrige}

% ===========================================================================
\section{Construire le logarithme népérien}

\begin{definition}[logarithme népérien]
La fonction $x \mapsto \dfrac{1}{x}$ est continue sur $\Rps$ : elle y admet
donc des primitives. On appelle \textbf{logarithme népérien}, et l'on note
$\ln$, l'unique primitive de $x \mapsto \dfrac{1}{x}$ sur $\Rps$ qui s'annule
en $1$. Autrement dit :
\[
  \ln'(x) = \frac{1}{x} \quad \text{pour tout } x>0,
  \qquad \text{et} \qquad \ln 1 = 0 .
\]
\end{definition}

\begin{propriete}[premières conséquences]
La fonction $\ln$ est définie et dérivable sur $\Rps$. Sa dérivée
$\dfrac{1}{x}$ y étant strictement positive, $\ln$ est \textbf{strictement
croissante} sur $\Rps$. De plus $\ln 1 = 0$, donc
\[
  \ln x > 0 \ \text{ si } x>1, \qquad \ln x < 0 \ \text{ si } 0<x<1 .
\]
\end{propriete}

\begin{theoreme}[propriétés algébriques]
Pour tous réels $a>0$ et $b>0$, tout entier $n$ et tout réel $x$ :
\[
  \ln(ab) = \ln a + \ln b, \qquad
  \ln\!\left(\frac{a}{b}\right) = \ln a - \ln b, \qquad
  \ln\!\left(\frac{1}{b}\right) = -\ln b,
\]
\[
  \ln\left(a^{n}\right) = n\ln a, \qquad
  \ln\sqrt{a} = \frac{1}{2}\ln a .
\]
\end{theoreme}

\begin{demonstration}[de la relation fondamentale]
Fixons $a>0$ et considérons la fonction $g$ définie sur $\Rps$ par
$g(x) = \ln(ax) - \ln x$. Elle est dérivable et
\[
  g'(x) = a \times \frac{1}{ax} - \frac{1}{x} = \frac{1}{x} - \frac{1}{x} = 0 .
\]
Sur l'intervalle $\Rps$, une fonction de dérivée nulle est constante : $g$ est
donc égale à sa valeur en $1$, soit $g(1) = \ln a - \ln 1 = \ln a$. Ainsi
$\ln(ax) - \ln x = \ln a$ pour tout $x>0$, ce qui est exactement
$\ln(ax) = \ln a + \ln x$. Les autres formules s'en déduisent : par exemple
$\ln b + \ln\frac1b = \ln 1 = 0$.
\end{demonstration}

\bmtmarge{Cette démonstration est un modèle du genre : pour établir une
identité, on introduit la différence des deux membres et l'on montre que sa
dérivée est nulle. Retenez le procédé, il resservira.}

\begin{propriete}[limites aux bornes]
\[
  \lim_{x \to 0^{+}} \ln x = -\infty, \qquad \limpinf \ln x = +\infty .
\]
La fonction $\ln$ réalise donc une bijection strictement croissante de $\Rps$
sur $\R$ : tout réel possède un antécédent, et un seul.
\end{propriete}

\begin{definition}[le nombre $e$]
Puisque $\ln$ est une bijection de $\Rps$ sur $\R$, il existe un unique réel
strictement positif dont le logarithme vaut $1$. On le note $e$ :
\[
  \ln e = 1, \qquad e \approx 2{,}718 .
\]
\end{definition}

\begin{definition}[la notation $e^{\alpha}$]
La bijectivité de $\ln$ permet d'aller plus loin. Pour tout réel $\alpha$, il
existe un unique réel strictement positif dont le logarithme vaut $\alpha$ :
on le note $e^{\alpha}$. Autrement dit, pour $x>0$ et $\alpha$ réel,
\[
  x = e^{\alpha} \iff \ln x = \alpha .
\]
\end{definition}

\begin{remarque}
Cette notation n'est pas un caprice : elle est cohérente avec les puissances
déjà connues. Puisque $\ln\!\left(e \times e\right) = 1+1 = 2$, le nombre
$e^{2}$ ainsi défini est bien le carré de $e$ ; le même raisonnement vaut pour
$e^{n}$, pour $e^{-n}$ et pour $e^{1/n}$. Le chapitre suivant reprendra
$\alpha \mapsto e^{\alpha}$ comme une fonction à part entière et en établira
les propriétés ; ici, elle nous sert seulement à écrire les solutions des
équations en $\ln$.
\end{remarque}

% ===========================================================================
\section{Limites usuelles et dérivées}

\begin{theoreme}[limites à connaître par cœur]
\[
  \limpinf \frac{\ln x}{x^{n}} = 0 \ (n \geq 1), \qquad
  \lim_{x \to 0^{+}} x\ln x = 0,
\]
\[
  \lim_{x \to 0} \frac{\ln(1+x)}{x} = 1, \qquad
  \lim_{x \to 1} \frac{\ln x}{x-1} = 1 .
\]
\end{theoreme}

\begin{remarque}
Les deux premières s'énoncent d'une phrase : \emph{en l'infini, le logarithme
est écrasé par toute puissance de $x$ ; en zéro, il est écrasé par $x$
lui-même.} Cette hiérarchie règle à elle seule la plupart des formes
indéterminées du chapitre. Les deux dernières ne sont rien d'autre que le
nombre dérivé de $\ln$ en $1$, écrit de deux façons.
\end{remarque}

\begin{theoreme}[dérivée d'un logarithme composé]
Soit $u$ une fonction dérivable et strictement positive sur un intervalle $I$.
Alors $\ln \circ\, u$ est dérivable sur $I$ et
\[
  \bigl(\ln u\bigr)' = \frac{u'}{u} .
\]
\end{theoreme}

\begin{theoreme}[la primitive qui manquait]
Soit $u$ dérivable et ne s'annulant pas sur un intervalle $I$. Alors
\[
  \int \frac{u'(x)}{u(x)}\dx = \ln\abs{u(x)} + k .
\]
En particulier, une primitive de $x \mapsto \dfrac{1}{x}$ sur
$\left]-\infty;0\right[$ est $x \mapsto \ln(-x)$.
\end{theoreme}

\begin{exemple}[reconnaître la forme $u'/u$]
Calculons $\displaystyle\int_{0}^{1} \frac{2x}{x^{2}+1}\dx$.

Ici $u(x) = x^{2}+1$, strictement positive, et $u'(x) = 2x$ : l'expression est
exactement $\dfrac{u'}{u}$. Donc
\[
  \int_{0}^{1} \frac{2x}{x^{2}+1}\dx = \Bigl[\ln\left(x^{2}+1\right)\Bigr]_{0}^{1}
  = \ln 2 - \ln 1 = \ln 2 .
\]

\smallskip
Si le numérateur n'est pas exactement $u'$, on ajuste par une constante :
$\displaystyle\int \frac{x}{x^{2}+1}\dx = \frac{1}{2}\ln\left(x^{2}+1\right)+k$.
\end{exemple}

% ===========================================================================
\section{L'étude complète de la fonction $\ln$}

\begin{exemple}[la courbe du logarithme]
\emph{Définition.} $\Df = \Rps$.

\emph{Limites.} $\lim_{x\to0^{+}} \ln x = -\infty$ : l'axe des ordonnées est
asymptote verticale. $\limpinf \ln x = +\infty$.

\emph{Branche infinie.} $\limpinf \dfrac{\ln x}{x} = 0$ : la courbe admet une
branche parabolique de direction $(Ox)$. Elle monte indéfiniment, mais de plus
en plus lentement.

\emph{Dérivée.} $\ln'(x) = \dfrac{1}{x} > 0$ : la fonction est strictement
croissante.

\emph{Convexité.} $\ln''(x) = -\dfrac{1}{x^{2}} < 0$ : la courbe est concave
sur tout $\Rps$, donc située sous chacune de ses tangentes.

\emph{Tangente en $1$.} $\ln 1 = 0$ et $\ln'(1) = 1$ : la tangente a pour
équation $y = x-1$. La concavité donne alors, pour tout $x>0$,
\[
  \ln x \leq x-1,
\]
inégalité classique qu'il est bon de savoir retrouver.
\end{exemple}

\begin{visualisation}[une croissance qui ne s'arrête jamais mais ralentit toujours]
\begin{center}
\begin{tikzpicture}
\begin{axis}[
  width = 11cm, height = 5.8cm,
  axis lines = middle, xlabel = {$x$},
  xmin = -0.6, xmax = 7.4, ymin = -3.2, ymax = 2.6,
  xtick = {1, 2.718}, xticklabels = {$1$, $e$}, ytick = {1},
  axis line style = {bmtGris},
  tick label style = {font=\footnotesize, color=bmtGris}, clip = true,
]
  \addplot[bmtIndigo, line width=1.3pt, domain=0.045:7.2, samples=250] {ln(x)};
  \addplot[bmtLaterite, dashed, line width=0.85pt, domain=0:3.4, samples=2] {x-1};
  \draw[bmtGris, dotted] (axis cs:2.718,0) -- (axis cs:2.718,1);
  \draw[bmtGris, dotted] (axis cs:0,1) -- (axis cs:2.718,1);
  \filldraw[bmtIndigo] (axis cs:1,0) circle (1.7pt);
  \node[bmtLaterite, font=\footnotesize, anchor=south west] at (axis cs:2.1,1.35) {$y=x-1$};
\end{axis}
\end{tikzpicture}
\end{center}

Deux choses à retenir de ce dessin. D'abord, la courbe reste toujours sous la
droite rouge, et ne la touche qu'en un seul point, celui d'abscisse $1$ :
c'est l'inégalité $\ln x \leq x-1$, lue au lieu d'être calculée. Ensuite,
regardez à quel point la montée s'aplatit vers la droite — il faut atteindre
$x = e \approx 2{,}7$ pour que le logarithme vaille seulement $1$, et
$x \approx 22\,000$ pour qu'il atteigne $10$.
\end{visualisation}

\begin{entrainement}[calculs et dérivées]
\exo[\niveau{1}] Simplifier $\ln 8 - \ln 2$, puis $\ln\sqrt{5} + \frac12\ln 5$.

\exo[\niveau{1}] Dériver $x \mapsto \ln(3x+1)$ sur $\left]-\frac13;+\infty\right[$,
puis $x \mapsto \ln\left(x^{2}+1\right)$ sur $\R$.

\exo[\niveau{2}] Déterminer $\limpinf \left[\ln(x+1) - \ln x\right]$.

\exo[\niveau{2}] Calculer $\displaystyle\int_{1}^{2}\frac{3x^{2}}{x^{3}+1}\dx$.

\exo[\niveau{3}] Étudier les variations de $f(x) = \dfrac{\ln x}{x}$ sur
$\Rps$ et déterminer son maximum.
\end{entrainement}

\begin{corrige}
\textbf{1.} $\ln 8 - \ln 2 = \ln 4 = 2\ln 2$ ; \
$\ln\sqrt5 + \frac12\ln5 = \frac12\ln5+\frac12\ln5 = \ln 5$.

\textbf{2.} $\dfrac{3}{3x+1}$ ; \ $\dfrac{2x}{x^{2}+1}$.

\textbf{3.} $\ln(x+1)-\ln x = \ln\dfrac{x+1}{x} = \ln\left(1+\dfrac1x\right) \to \ln 1 = 0$.

\textbf{4.} C'est $\dfrac{u'}{u}$ avec $u = x^{3}+1$ :
$\bigl[\ln(x^{3}+1)\bigr]_{1}^{2} = \ln 9 - \ln 2 = \ln\dfrac92$.

\textbf{5.} $f'(x) = \dfrac{1-\ln x}{x^{2}}$, du signe de $1-\ln x$ : $f$ croît
sur $\intervalleof{0}{e}$, décroît ensuite. Maximum
$f(e) = \dfrac{1}{e}$.
\end{corrige}

% ===========================================================================
\section{Équations, inéquations et systèmes}

\begin{methode}{résoudre une équation avec des logarithmes}
Trois étapes, dans cet ordre, et la première n'est jamais facultative.

\begin{enumerate}[label=\textbf{\arabic*.}, leftmargin=*, itemsep=0.3em]
  \item \textbf{Conditions d'existence.} Écrire que chaque expression placée
        sous un $\ln$ est strictement positive, et résoudre ce système. On
        obtient un ensemble $D$ hors duquel aucune solution n'est recevable.
  \item \textbf{Résolution.} Utiliser la stricte croissance de $\ln$, qui rend
        la fonction injective :
        \[
          \ln A = \ln B \daccord A = B, \qquad
          \ln A \leq \ln B \daccord A \leq B \quad (A,B>0).
        \]
  \item \textbf{Vérification.} Ne conserver que les solutions appartenant
        à $D$. C'est ici que se perdent le plus de points.
\end{enumerate}
\end{methode}

\begin{exemple}[une équation où une solution doit être écartée]
Résolvons $\ln(x-1) + \ln(x+2) = \ln 4$.

\emph{Conditions d'existence.} $x-1>0$ et $x+2>0$, soit $x>1$. Donc
$D = \left]1;+\infty\right[$.

\emph{Résolution.} Le membre de gauche vaut $\ln\bigl[(x-1)(x+2)\bigr]$.
L'équation devient $(x-1)(x+2) = 4$, soit $x^{2}+x-6 = 0$, dont les racines
sont $2$ et $-3$.

\emph{Vérification.} Seul $2$ appartient à $D$. L'ensemble des solutions est
donc $\{2\}$ — la valeur $-3$ est à rejeter, car $\ln(-4)$ n'a aucun sens.
\end{exemple}

\begin{erreurclassique}
Passer directement de $\ln A + \ln B = \ln C$ à $AB = C$ sans avoir posé les
conditions d'existence donne parfois la bonne réponse, mais fait toujours
perdre des points — et, comme ci-dessus, conduit régulièrement à retenir une
solution qui n'en est pas une. Posez les conditions \emph{avant} de calculer,
pas après.
\end{erreurclassique}

\begin{exemple}[un système logarithmique]
Résolvons le système
$\begin{cases} \ln x + \ln y = \ln 12 \\ \ln x - \ln y = \ln 3 \end{cases}$
d'inconnues $x>0$ et $y>0$.

En additionnant les deux lignes : $2\ln x = \ln 36$, donc $\ln x = \ln 6$ et
$x = 6$. En soustrayant : $2\ln y = \ln 12 - \ln 3 = \ln 4$, donc $y = 2$.

Les deux valeurs sont strictement positives : le couple $(6\,;2)$ est la
solution, et l'on vérifie que $\ln 6 + \ln 2 = \ln 12$.
\end{exemple}

\begin{entrainement}[équations et inéquations]
\exo[\niveau{1}] Résoudre $\ln(2x-3) = 0$, puis $\ln(x^{2}) = \ln(3x-2)$.

\exo[\niveau{2}] Résoudre l'inéquation $\ln(x-1) \leq \ln(5-x)$.

\exo[\niveau{2}] Résoudre $\ln(x+3) + \ln(x-1) = \ln 5$.

\exo[\niveau{3}] Résoudre le système
$\begin{cases} \ln x + 2\ln y = 5\ln 2 \\ 3\ln x - \ln y = \ln 2 \end{cases}$.
\end{entrainement}

\begin{corrige}
\textbf{1.} $2x-3 = 1$ avec $2x-3>0$ : $x = 2$. Pour la seconde, conditions
$x \neq 0$ et $3x-2>0$ ; l'équation $x^{2} = 3x-2$ donne $x=1$ ou $x=2$, tous
deux acceptables.

\textbf{2.} Conditions : $1<x<5$. L'inéquation équivaut à $x-1 \leq 5-x$, soit
$x \leq 3$. Solutions : $\intervalleof{1}{3}$.

\textbf{3.} Conditions : $x>1$. Puis $(x+3)(x-1) = 5$, soit $x^{2}+2x-8=0$,
racines $2$ et $-4$. Seul $2$ convient.

\textbf{4.} Posons $X = \ln x$ et $Y = \ln y$ : le système devient
$X+2Y = 5\ln2$ et $3X-Y = \ln2$. On trouve $X = \ln 2$ et $Y = 2\ln 2$, d'où
$x = 2$ et $y = 4$.
\end{corrige}

% ===========================================================================
\section{Logarithme décimal et logarithme de base quelconque}

\begin{definition}[logarithme de base $a$]
Soit $a>0$ avec $a \neq 1$. On appelle \textbf{logarithme de base $a$} la
fonction définie sur $\Rps$ par
\[
  \log_{a} x = \frac{\ln x}{\ln a}.
\]
Le cas $a = 10$ donne le \textbf{logarithme décimal}, noté $\log$ :
$\log x = \dfrac{\ln x}{\ln 10}$.
\end{definition}

\begin{propriete}
Le logarithme de base $a$ possède les mêmes propriétés algébriques que $\ln$.
Il est strictement croissant si $a>1$, strictement décroissant si $0<a<1$. De
plus $\log_{a} a = 1$ et $\log_{a} 1 = 0$.
\end{propriete}

\begin{remarque}[à quoi sert le logarithme décimal]
Pour un entier $N \geq 1$, le nombre de chiffres de son écriture décimale vaut
$\ent{\log N}+1$. Le logarithme décimal mesure donc, littéralement, l'ordre de
grandeur : $\log 1000 = 3$, $\log 10^{6} = 6$. C'est pour cela qu'il sert
d'échelle en chimie (le pH) et en acoustique (le décibel).
\end{remarque}

\begin{entrainement}[bases quelconques]
\exo[\niveau{1}] Calculer $\log 100$, $\log_{2} 8$ et $\log_{3} \frac{1}{9}$.

\exo[\niveau{2}] Combien de chiffres compte $2^{100}$ ? On donne
$\log 2 \approx 0{,}301$.

\exo[\niveau{3}] Résoudre $\log_{2} x + \log_{4} x = 3$. \emph{On exprimera
tout à l'aide de $\ln$.}
\end{entrainement}

\begin{corrige}
\textbf{1.} $2$ ; \ $3$ ; \ $-2$.

\textbf{2.} $\log\left(2^{100}\right) = 100\log 2 \approx 30{,}1$, donc
$\ent{30{,}1}+1 = 31$ chiffres.

\textbf{3.} $\dfrac{\ln x}{\ln 2} + \dfrac{\ln x}{2\ln 2} = 3$, soit
$\dfrac{3\ln x}{2\ln 2} = 3$ et $\ln x = 2\ln 2$ : $x = 4$.
\end{corrige}

% ===========================================================================
% ===========================================================================
\begin{annales}
Le logarithme est présent dans presque tous les sujets, et rarement seul : il
vient accompagné d'une fonction auxiliaire, d'un encadrement ou d'une
bijection. Les énoncés qui suivent sont repris de sessions réelles ; leur
provenance est indiquée à chaque fois. Plusieurs d'entre eux étaient, dans le
sujet d'origine, des questions de limite, de dérivation ou de calcul
intégral : la notion qu'ils font travailler en plus du logarithme est signalée
entre parenthèses, à la suite de la provenance.

\exoann[limites et continuité]{Baccalauréat, Madagascar, série D, session 2026 — problème}
Soit $f(x) = x^{2}(1+\ln x)+x-1$ sur $\intervalleo{0}{+\infty}$. Calculer
$\lim_{x\to0^{+}}f(x)$, puis étudier le signe de $f'(x)$ et dresser le
tableau de variation.

\exoann[limites et continuité]{Baccalauréat, Madagascar, série C, session 2026 — problème II, partie A}
Pour $n\in\N^{*}$, soit $f_{n}(x) = x^{n}\ln(x+1)$, définie sur
$\intervalleo{-1}{+\infty}$. Étudier, suivant la parité de $n$, la limite de
$f_{n}(x)$ à droite de $-1$. Qu'en conclut-on pour la courbe ?

\exoann[dérivabilité]{Baccalauréat, Congo-Brazzaville, série C, session 2023 — exercice 3}
Soit $f(x) = \dfrac{6\ln(x+3)}{x+3}$ sur $\intervallefo{0}{+\infty}$.
Démontrer que $f'(x) = \dfrac{6\left[1-\ln(x+3)\right]}{(x+3)^{2}}$, en
déduire le sens de variation de $f$, puis calculer
$\lim_{x\to+\infty}f(x)$ et dresser le tableau de variation.

\exoann[dérivabilité]{Baccalauréat, Madagascar, série C, session 2026 — problème II}
Pour $n\in\N^{*}$, soit $f_{n}(x) = x^{n}\ln(x+1)$ sur
$\intervalleo{-1}{+\infty}$. Calculer $f_{n}'(x)$, puis montrer que
$f_{n}'(x)>0$ pour tout $x>0$.

\exoann[dérivabilité]{Baccalauréat, Cameroun, série C, session 2023 — exercice 1}
Sur $\intervalleo{1}{+\infty}$ on pose $f(x) = x-2+\ln\sqrt{x-1}$ et
$h(x) = \dfrac{2x-2-\ln(x-1)}{2\sqrt{x-1}}$. Montrer que
$h'(x) = \dfrac{f(x)}{2(x-1)\sqrt{x-1}}$, puis en déduire les variations de
$h$, sachant que $f$ est strictement croissante et s'annule en $2$.

\exoann[étude de fonction]{Baccalauréat, Mauritanie, série C, session 2022 — exercice 4}
Sur $\intervalleo{0}{+\infty}$ on pose $f(x) = x-3-\dfrac{3\ln x}{x}$ et
$u(x) = x^{2}-3+3\ln x$.
\begin{enumerate}[label=\textbf{\arabic*)}, leftmargin=*, itemsep=0.15em]
  \item Dresser le tableau de variation de $u$, montrer que $u(x) = 0$ admet
        une solution unique $\alpha$ avec $1{,}4 < \alpha < 1{,}41$, puis
        donner le signe de $u$.
  \item Calculer les limites de $f$ en $0^{+}$ et en $+\infty$, montrer que la
        droite $(D) : y = x-3$ est asymptote et préciser la position relative.
  \item Sachant que $f'(x) = \dfrac{u(x)}{x^{2}}$, dresser le tableau de
        variation de $f$ et construire sa courbe.
  \item Soit $\varphi$ la restriction de $f$ à $\intervalleof{0}{\alpha}$.
        Montrer que $\varphi$ est une bijection de $\intervalleof{0}{\alpha}$
        sur un intervalle $J$ à déterminer, puis construire la courbe de
        $\varphi^{-1}$.
\end{enumerate}

\exoann[étude de fonction]{Baccalauréat, Burkina Faso, série D, session 2023 — problème, parties A et B}
Sur $\left]-\infty\,;0\right[$ on pose $h(x) = 2\ln(-x)$, et $f$ est définie
par $f(x) = -2x+1+2x\ln\abs{x}$ pour $x<0$.
\begin{enumerate}[label=\textbf{\alph*)}, leftmargin=*, itemsep=0.15em]
  \item Étudier les variations de $h$ et dresser son tableau.
  \item Déterminer le signe de $h(x)$ suivant les valeurs de $x$.
  \item Calculer $f'(x)$ pour $x<0$ et reconnaître $h$. En déduire le sens de
        variation de $f$ sur $\left]-\infty\,;0\right[$.
\end{enumerate}

\exoann[étude de fonction]{Baccalauréat, Cameroun, série C, session 2023 — exercice 1, suite}
Sur $\intervalleo{1}{+\infty}$ on pose $f(x) = x-2+\ln\sqrt{x-1}$.
\begin{enumerate}[label=\textbf{\alph*)}, leftmargin=*, itemsep=0.15em]
  \item Dresser le tableau de variation de $f$.
  \item Montrer que $2$ est l'unique solution de $f(x) = 0$.
  \item En déduire le signe de $f(x)$ suivant les valeurs de $x$.
\end{enumerate}

\exoann{Baccalauréat, Côte d'Ivoire, série D, session 2023 — exercice 2}
Résoudre dans $\R$ l'inéquation $\ln(1-x) < 2$.

\exoann[calcul intégral]{Baccalauréat, Congo-Brazzaville, série C, session 2023 — exercice 3, suite}
On garde $f(x) = \dfrac{6\ln(x+3)}{x+3}$ et l'on pose
$g(x) = \ln^{2}(x+3)$,
$u_{n} = \displaystyle\int_{n}^{n+1}f(x)\dx$ et
$I_{n} = \displaystyle\int_{0}^{n}f(x)\dx$.
\begin{enumerate}[label=\textbf{\alph*)}, leftmargin=*, itemsep=0.15em]
  \item Sachant que $f$ est décroissante sur $\intervallefo{0}{+\infty}$,
        établir $f(n+1) \leq u_{n} \leq f(n)$, puis calculer
        $\lim_{n\to+\infty}u_{n}$.
  \item Montrer que $g'(x) = \frac13 f(x)$ et en déduire $I_{n}$.
\end{enumerate}

\exoann[calcul intégral]{Baccalauréat, Cameroun, série C, session 2023 — exercice 1, fin}
On garde $h(x) = \dfrac{2x-2-\ln(x-1)}{2\sqrt{x-1}}$. Calculer
$\displaystyle\int_{2}^{3}\frac{\ln(x-1)}{2\sqrt{x-1}}\dx$, puis en déduire
$I = \displaystyle\int_{2}^{3}h(x)\dx$.

\exoann[calcul intégral]{Baccalauréat, Mauritanie, série C, session 2022 — exercice 4, fin}
On garde $f(x) = x-3-\dfrac{3\ln x}{x}$. Calculer l'aire, en unités d'aire,
de la partie du plan délimitée par la courbe de $f$, l'axe des abscisses et
les droites d'équations $x = 1$ et $x = e$.

\exoann[calcul intégral]{Baccalauréat, Madagascar, série D, session 2026 — problème, fin}
Calculer $\displaystyle\int_{1}^{e}x^{2}\ln x\dx$ par une intégration par
parties.

\exoann[calcul intégral]{Baccalauréat, Madagascar, série C, session 2026 — problème II, fin}
Pour $n\in\N^{*}$, on pose
$I_{n} = \displaystyle\int_{0}^{1}t^{n}\ln(t+1)\dt$. Calculer $I_{1}$ par une
intégration par parties, puis montrer que $(I_{n})$ est décroissante et
converge vers $0$.
\end{annales}

\begin{corrige}[annales]
\textbf{1.} En $0^{+}$, $x^{2}\ln x\to0$ et $x^{2}\to0$, donc $f(x)\to-1$. On
calcule $f'(x) = 3x+2x\ln x+1$. En posant $\varphi = f'$, on a
$\varphi'(x) = 5+2\ln x$, qui s'annule en $e^{-5/2}$ où
$\varphi\!\left(e^{-5/2}\right) = 1-2e^{-5/2}>0$ : le minimum de $f'$ est
strictement positif, donc $f$ est strictement croissante sur
$\intervalleo{0}{+\infty}$, de $-1$ à $+\infty$.

\textbf{2.} Quand $x\to-1^{+}$, $\ln(x+1)\to-\infty$ et $x^{n}\to(-1)^{n}$. Si
$n$ est pair, $f_{n}(x)\to-\infty$ ; si $n$ est impair, $f_{n}(x)\to+\infty$.
Dans les deux cas la droite $x = -1$ est asymptote verticale à la courbe.

\textbf{3.} $f'(x) = 6\,\dfrac{\frac{1}{x+3}(x+3)-\ln(x+3)}{(x+3)^{2}}
= \dfrac{6\left[1-\ln(x+3)\right]}{(x+3)^{2}}$. Pour $x \geq 0$ on a
$x+3 \geq 3 > e$, donc $\ln(x+3)>1$ et $f'(x)<0$ : $f$ est strictement
décroissante. Comme $\frac{\ln X}{X}\to0$, $f(x)\to0$ en $+\infty$ ; $f$
décroît donc de $f(0) = 2\ln3$ à $0$.

\textbf{4.} $f_{n}'(x) = nx^{n-1}\ln(x+1)+\dfrac{x^{n}}{x+1}$. Pour $x>0$ les
trois quantités $x^{n-1}$, $\ln(x+1)$ et $\frac{x^{n}}{x+1}$ sont strictement
positives : la somme l'est aussi.

\textbf{5.} En posant $t = x-1>0$, $h = \sqrt t-\dfrac{\ln t}{2\sqrt t}$ et le
calcul donne $h' = \dfrac{t-1+\frac12\ln t}{2t\sqrt t}
= \dfrac{f(x)}{2(x-1)\sqrt{x-1}}$. Le dénominateur étant positif, $h'$ a le
signe de $f$ : $h$ décroît sur $\intervalleof{1}{2}$, croît sur
$\intervallefo{2}{+\infty}$, et atteint en $2$ son minimum $h(2) = 1$.

\textbf{6.} $u'(x) = 2x+\frac3x>0$ : $u$ croît de $-\infty$ à $+\infty$, donc
s'annule une fois et une seule. Comme $u(1{,}4)<0<u(1{,}41)$, la racine
$\alpha$ est dans $\intervalleo{1{,}4}{1{,}41}$ ; $u$ est négative avant
$\alpha$, positive après. En $0^{+}$, $\frac{\ln x}{x}\to-\infty$ donc
$f(x)\to+\infty$ ; en $+\infty$, $\frac{\ln x}{x}\to0$ et
$f(x)-(x-3) = -\frac{3\ln x}{x}\to0$ : la droite $y = x-3$ est asymptote, la
courbe est au-dessus pour $x<1$ et en dessous pour $x>1$. Le signe de $f'$
étant celui de $u$, $f$ décroît sur $\intervalleof{0}{\alpha}$ et croît
ensuite. Enfin $\varphi$ est continue et strictement décroissante de
$\intervalleof{0}{\alpha}$ sur $J = \intervallefo{f(\alpha)}{+\infty}$ :
c'est une bijection, et la courbe de $\varphi^{-1}$ s'obtient par symétrie
par rapport à la droite $y = x$.

\textbf{7. a)} $h'(x) = \frac2x<0$ sur $\left]-\infty\,;0\right[$ : $h$ décroît
de $+\infty$ à $-\infty$.
\textbf{b)} $h(x) = 0$ équivaut à $-x = 1$, soit $x = -1$ : $h$ est positive
sur $\left]-\infty\,;-1\right[$, négative sur $\intervalleo{-1}{0}$.
\textbf{c)} $f'(x) = -2+2\ln\abs{x}+2 = 2\ln(-x) = h(x)$. Donc $f$ croît sur
$\left]-\infty\,;-1\right]$ et décroît sur $\intervalleo{-1}{0}$.

\textbf{8. a)} $f(x) = x-2+\frac12\ln(x-1)$ et
$f'(x) = 1+\frac{1}{2(x-1)}>0$ : $f$ croît de $-\infty$ à $+\infty$.
\textbf{b)} La stricte croissance donne au plus une racine, et $f(2) = 0$.
\textbf{c)} $f$ est négative sur $\intervalleo{1}{2}$, positive sur
$\intervalleo{2}{+\infty}$.

\textbf{9.} L'existence impose $1-x>0$, soit $x<1$. L'inéquation équivaut
alors à $1-x<e^{2}$, soit $x>1-e^{2}$. L'ensemble des solutions est
$\intervalleo{1-e^{2}}{1}$.

\textbf{10. a)} Sur $\intervalle{n}{n+1}$, la décroissance de $f$ donne
$f(n+1) \leq f(x) \leq f(n)$ ; en intégrant sur un intervalle de longueur
$1$, $f(n+1) \leq u_{n} \leq f(n)$. Les deux bornes tendent vers $0$, donc
$u_{n}\to0$.
\textbf{b)} $g'(x) = \dfrac{2\ln(x+3)}{x+3} = \dfrac13 f(x)$, donc $3g$ est une
primitive de $f$ et
$I_{n} = 3\left[\ln^{2}(x+3)\right]_{0}^{n} = 3\left(\ln^{2}(n+3)-\ln^{2}3\right)$.

\textbf{11.} On pose $t = x-1$, puis on intègre par parties avec $u = \ln t$
et $v' = \frac{1}{2\sqrt t}$, donc $v = \sqrt t$ :
\[
  \int_{2}^{3}\frac{\ln(x-1)}{2\sqrt{x-1}}\dx
  = \left[\sqrt t\,\ln t\right]_{1}^{2}-\int_{1}^{2}\frac{\dt}{\sqrt t}
  = \sqrt2\ln2-2\sqrt2+2 .
\]
Comme $h(x) = \sqrt{x-1}-\dfrac{\ln(x-1)}{2\sqrt{x-1}}$ et
$\displaystyle\int_{2}^{3}\sqrt{x-1}\dx = \frac{4\sqrt2}{3}-\frac23$, il
vient $I = \dfrac{10\sqrt2}{3}-\sqrt2\ln2-\dfrac83$.

\textbf{12.} Sur $\intervalle{1}{e}$ la fonction $f$ est négative : l'aire
cherchée vaut $\displaystyle\int_{1}^{e}\left(3-x+\frac{3\ln x}{x}\right)\dx
= \left[3x-\frac{x^{2}}{2}+\frac32\ln^{2}x\right]_{1}^{e}
= 3e-\frac{e^{2}}{2}-1$ unités d'aire.

\textbf{13.} Avec $u = \ln x$ et $v' = x^{2}$ :
$\displaystyle\int_{1}^{e}x^{2}\ln x\dx
= \left[\frac{x^{3}}{3}\ln x\right]_{1}^{e}-\frac13\int_{1}^{e}x^{2}\dx
= \frac{e^{3}}{3}-\frac{e^{3}-1}{9} = \frac{2e^{3}+1}{9}$.

\textbf{14.} Avec $u = \ln(t+1)$ et $v' = t$ :
\[
  I_{1} = \left[\frac{t^{2}}{2}\ln(t+1)\right]_{0}^{1}
  -\frac12\int_{0}^{1}\frac{t^{2}}{t+1}\dt
  = \frac{\ln2}{2}-\frac12\left(-\frac12+\ln2\right) = \frac14 .
\]
Sur $\intervalle{0}{1}$ on a $t^{n+1} \leq t^{n}$ et $\ln(t+1)\geq0$, donc
$I_{n+1} \leq I_{n}$ : la suite décroît. Enfin
$0 \leq I_{n} \leq \ln2\displaystyle\int_{0}^{1}t^{n}\dt = \frac{\ln2}{n+1}$,
qui tend vers $0$ : $(I_{n})$ converge vers $0$.
\end{corrige}

\begin{jecherche}[la datation d'un charbon de bois]
Un site archéologique des Hautes Terres a livré un charbon de bois. La
proportion de carbone 14 restante dans un échantillon, $t$ années après la
mort de l'arbre, est modélisée par
\[
  p(t) = p_{0}\,\alpha^{t}, \qquad \text{avec } \alpha \in \intervalleo{0}{1},
\]
où $p_{0}$ est la proportion initiale. On sait que la moitié du carbone 14 a
disparu au bout de $5\,730$ ans.

\begin{enumerate}[label=\textbf{\arabic*.}, leftmargin=*, itemsep=0.3em]
  \item Traduire l'information sur la demi-vie par une équation vérifiée
        par $\alpha$, et exprimer $\ln \alpha$ en fonction de $\ln 2$.
  \item L'échantillon analysé contient $62\,\%$ de la proportion initiale.
        Exprimer l'âge $t$ du charbon en fonction de $\ln$, puis en donner une
        valeur approchée.
  \item Montrer que la fonction $t \mapsto p(t)$ est strictement décroissante,
        et déterminer sa limite en $+\infty$. Le modèle prévoit-il une
        disparition totale du carbone 14 ?
  \item Au bout de combien d'années ne reste-t-il plus que $1\,\%$ de la
        proportion initiale ?
  \item La méthode devient peu fiable en dessous de $1\,\%$. Quelle est, selon
        ce modèle, la portée maximale de la datation au carbone 14 ?
\end{enumerate}
\end{jecherche}

\begin{corrige}[la datation d'un charbon de bois]
\textbf{1.} La demi-vie s'écrit $p(5730) = \tfrac12 p_{0}$, soit
$\alpha^{5730} = \tfrac12$. En composant par le logarithme :
$5730\ln\alpha = -\ln 2$, donc
$\ln\alpha = -\dfrac{\ln 2}{5730}$, qui est bien négatif.

\textbf{2.} La condition s'écrit $\alpha^{t} = 0{,}62$, donc
$t\ln\alpha = \ln 0{,}62$ et
\[
  t = \frac{\ln 0{,}62}{\ln\alpha} = -\,\frac{5730\,\ln 0{,}62}{\ln 2}
  \approx 3\,950 \text{ ans}.
\]

\textbf{3.} $p(t) = p_{0}e^{t\ln\alpha}$ avec $\ln\alpha < 0$ : la fonction est
strictement décroissante et tend vers $0$ en $+\infty$. Le modèle ne prévoit
donc \emph{pas} de disparition totale : la proportion s'approche de zéro sans
jamais l'atteindre en temps fini.

\textbf{4.} $\alpha^{t} = 0{,}01$ donne
$t = -\dfrac{5730\ln 0{,}01}{\ln 2} \approx 38\,000$ ans.

\textbf{5.} La méthode porte donc jusqu'à quarante mille ans environ. Au-delà,
la quantité résiduelle est trop faible pour être mesurée de façon fiable —
c'est une limite de l'instrument autant que du modèle.
\end{corrige}


% ===========================================================================
\begin{jemeteste}
Répondre par vrai ou faux, en justifiant en une ligne.

\begin{enumerate}[label=\textbf{\arabic*.}, leftmargin=*, itemsep=0.28em]
  \item $\ln(a+b) = \ln a + \ln b$ pour tous $a>0$ et $b>0$.
  \item L'équation $\ln x = -5$ n'a pas de solution.
  \item Pour tout $x>0$, $\ln x \leq x-1$.
  \item Une primitive de $x \mapsto \dfrac{1}{x}$ sur $\left]-\infty;0\right[$
        est $x \mapsto \ln x$.
  \item $\limpinf \dfrac{\ln x}{x^{2}} = 0$.
  \item L'inéquation $\ln(x-2) \leq 3$ a pour solutions $x \leq e^{3}+2$.
\end{enumerate}
\end{jemeteste}

\begin{corrige}[je me teste]
\textbf{1.} Faux — c'est $\ln(ab)$ qui vaut $\ln a + \ln b$.
\textbf{2.} Faux — $\ln$ est une bijection de $\Rps$ sur $\R$, donc tout réel
est atteint ; ici $x = e^{-5}$.
\textbf{3.} Vrai — concavité et tangente en $1$.
\textbf{4.} Faux — $\ln x$ n'y est pas défini ; c'est $\ln(-x)$.
\textbf{5.} Vrai.
\textbf{6.} Faux — il manque la condition d'existence $x>2$ : les solutions
sont $\intervalleof{2}{e^{3}+2}$.
\end{corrige}

% ===========================================================================
\begin{commeaubac}
\bmtsource{Baccalauréat, session 2026 — série C, problème II, partie A}
Soit $f_{n}$ la fonction définie sur $\left]-1;+\infty\right[$ par
$f_{n}(x) = x^{n}\ln(x+1)$, pour tout $n \in \N^{*}$. On note $(C_{n})$ ses
courbes représentatives dans un repère orthonormé d'unité $2$ cm.

\begin{enumerate}[label=\textbf{\arabic*.}, leftmargin=*, itemsep=0.25em]
  \item \begin{enumerate}[label=\textbf{\alph*)}, leftmargin=1.4em]
          \item Étudier, suivant la parité de $n$, la limite de $f_{n}(x)$ à
                droite de $-1$. Que peut-on conclure ?
                \hfill\textup{(0,75 pt)}
          \item Calculer $\limpinf f_{n}(x)$. \hfill\textup{(0,25 pt)}
        \end{enumerate}
  \item Montrer que pour tout $x \in \left]-1;+\infty\right[$ et tout
        $n \in \N^{*}$, les courbes $(C_{n})$ passent par deux points fixes
        dont on déterminera les coordonnées. \hfill\textup{(0,5 pt)}
  \item Soit $g_{n}$ définie sur $\left]-1;+\infty\right[$ par
        $g_{n}(x) = n\ln(x+1) + \dfrac{x}{x+1}$.
        \begin{enumerate}[label=\textbf{\alph*)}, leftmargin=1.4em]
          \item Dresser le tableau de variation de $g_{n}$ sur
                $\left]-1;+\infty\right[$. \hfill\textup{(0,75 pt)}
          \item Calculer $g_{n}(0)$ et en déduire le signe de $g_{n}(x)$
                suivant les valeurs de $x$. \hfill\textup{(0,5 pt)}
        \end{enumerate}
  \item \begin{enumerate}[label=\textbf{\alph*)}, leftmargin=1.4em]
          \item Montrer que $f'_{n}(x) = x^{n-1}g_{n}(x)$.
                \hfill\textup{(0,5 pt)}
          \item Dresser le tableau de variation de $f_{n}$ selon la parité
                de $n$. \hfill\textup{(1 pt)}
        \end{enumerate}
  \item Construire les courbes $(C_{1})$ et $(C_{2})$ dans un même repère.
        \hfill\textup{(1,25 pt)}
\end{enumerate}
\end{commeaubac}

\begin{corrige}[baccalauréat 2026, série C]
\textbf{1. a)} Quand $x \to -1^{+}$, $\ln(x+1) \to -\infty$ et
$x^{n} \to (-1)^{n}$. Si $n$ est pair, $x^{n} \to 1$ et
$f_{n}(x) \to -\infty$ ; si $n$ est impair, $x^{n} \to -1$ et
$f_{n}(x) \to +\infty$. Dans les deux cas la droite $x = -1$ est asymptote
verticale. \textbf{b)} En $+\infty$, les deux facteurs tendent vers $+\infty$ :
$f_{n}(x) \to +\infty$.

\textbf{2.} $f_{n}(x)$ ne dépend pas de $n$ lorsque $x^{n}$ est indépendant de
$n$, c'est-à-dire pour $x = 0$ et $x = 1$ ; ou bien lorsque le second facteur
est nul, ce qui donne aussi $x = 0$. Les points fixes sont $(0;0)$ et
$(1;\ln 2)$.

\textbf{3. a)} $g'_{n}(x) = \dfrac{n}{x+1} + \dfrac{1}{(x+1)^{2}}
= \dfrac{n(x+1)+1}{(x+1)^{2}}$, strictement positive sur
$\left]-1;+\infty\right[$ dès que $n \geq 1$ : $g_{n}$ est strictement
croissante, de $-\infty$ à $+\infty$. \textbf{b)} $g_{n}(0) = 0$ : donc
$g_{n}(x) < 0$ pour $x<0$ et $g_{n}(x)>0$ pour $x>0$.

\textbf{4. a)} $f'_{n}(x) = nx^{n-1}\ln(x+1) + \dfrac{x^{n}}{x+1}
= x^{n-1}\left[n\ln(x+1) + \dfrac{x}{x+1}\right] = x^{n-1}g_{n}(x)$.
\textbf{b)} Le signe de $f'_{n}$ est celui de $x^{n-1}g_{n}(x)$. Si $n$ est
impair, $x^{n-1} \geq 0$ et $f'_{n}$ a le signe de $g_{n}$ : $f_{n}$ décroît
sur $\intervalleo{-1}{0}$ puis croît. Si $n$ est pair, $x^{n-1}$ a le signe de
$x$, donc $f'_{n} = x^{n-1}g_{n} \geq 0$ partout : $f_{n}$ est croissante.

\textbf{5.} Les deux courbes passent par $(0;0)$ et $(1;\ln2)$, admettent
$x=-1$ pour asymptote verticale, et se distinguent par leur comportement au
voisinage de $-1$ (branche vers $+\infty$ pour $(C_{1})$, vers $-\infty$ pour
$(C_{2})$).
\end{corrige}

\begin{commeaubac}
\bmtsource{Exercice au format de l'épreuve}
Soit $f$ la fonction définie sur $\Rps$ par
$f(x) = x - 2 + \dfrac{\ln x}{x}$, de courbe $(\Cf)$.

\begin{enumerate}[label=\textbf{\arabic*.}, leftmargin=*, itemsep=0.25em]
  \item Déterminer les limites de $f$ en $0^{+}$ et en $+\infty$.
        \hfill\textup{(1 pt)}
  \item Montrer que la droite $(\mathscr{D})$ d'équation $y = x-2$ est
        asymptote à $(\Cf)$ en $+\infty$, et étudier la position relative de la
        courbe et de la droite. \hfill\textup{(1,25 pt)}
  \item Montrer que $f'(x) = \dfrac{x^{2}+1-\ln x}{x^{2}}$, puis établir que
        $f'(x)>0$ pour tout $x>0$. \emph{On pourra utiliser l'inégalité
        $\ln x \leq x-1$.} \hfill\textup{(1,5 pt)}
  \item Dresser le tableau de variation de $f$. \hfill\textup{(0,75 pt)}
  \item Montrer que l'équation $f(x)=0$ admet une unique solution $\alpha$ et
        que $1 < \alpha < 2$. \hfill\textup{(1 pt)}
  \item Calculer $\displaystyle\int_{1}^{e} \frac{\ln x}{x}\dx$, puis l'aire du
        domaine délimité par $(\Cf)$, la droite $(\mathscr{D})$ et les droites
        $x=1$ et $x=e$. \hfill\textup{(1,5 pt)}
\end{enumerate}
\end{commeaubac}

\begin{corrige}[exercice au format de l'épreuve]
\textbf{1.} En $0^{+}$ : $x-2 \to -2$ et $\dfrac{\ln x}{x} \to -\infty$
(numérateur $\to-\infty$, dénominateur $\to 0^{+}$), donc $f(x) \to -\infty$.
En $+\infty$ : $\dfrac{\ln x}{x} \to 0$ donc $f(x) \to +\infty$.

\textbf{2.} $f(x)-(x-2) = \dfrac{\ln x}{x} \to 0$ : asymptote. Le signe de
cette différence est celui de $\ln x$ : la courbe est au-dessus de
$(\mathscr{D})$ pour $x>1$, au-dessous pour $0<x<1$, et les deux se coupent
en $x=1$.

\textbf{3.} $f'(x) = 1 + \dfrac{1-\ln x}{x^{2}} = \dfrac{x^{2}+1-\ln x}{x^{2}}$.
Comme $\ln x \leq x-1 \leq x^{2}$ pour $x>0$, le numérateur est minoré par
$x^{2}+1-x^{2} = 1 > 0$ : donc $f'>0$.

\textbf{4.} $f$ est strictement croissante sur $\Rps$, de $-\infty$ à $+\infty$.

\textbf{5.} $f$ étant continue et strictement croissante, avec $f(1) = -1 < 0$
et $f(2) = \dfrac{\ln 2}{2} > 0$, le théorème des valeurs intermédiaires donne
une unique solution $\alpha \in \intervalleo{1}{2}$.

\textbf{6.} L'expression $\dfrac{\ln x}{x}$ est de la forme $u'u$ avec
$u = \ln x$ : une primitive est $\dfrac{(\ln x)^{2}}{2}$, d'où
$\displaystyle\int_{1}^{e}\frac{\ln x}{x}\dx = \frac12$. Sur
$\intervalle{1}{e}$ la courbe est au-dessus de la droite, donc l'aire vaut
$\dfrac12$ unité d'aire.
\end{corrige}

% ===========================================================================
\begin{concours}
\bmtsource{D'après le concours d'entrée de l'ENI, Fianarantsoa}
Résoudre dans $\Rps \times \Rps$ le système
\[
  \begin{cases}
    \dfrac{1}{\ln x} + \dfrac{1}{\ln y} = \dfrac{7}{3} \\[0.8em]
    \ln(xy) = \dfrac{7}{2}
  \end{cases}
\]

Puis calculer l'intégrale $\displaystyle J = \int_{1}^{3} \frac{(x+3)\dx}{(x+1)^{3}}$.
\end{concours}

\begin{corrige}[vers le concours]
\textbf{Le système.} Posons $X = \ln x$ et $Y = \ln y$, tous deux non nuls. La
seconde équation donne $X+Y = \dfrac{7}{2}$. La première s'écrit
$\dfrac{X+Y}{XY} = \dfrac{7}{3}$, donc
$\dfrac{7/2}{XY} = \dfrac{7}{3}$ et $XY = \dfrac{3}{2}$.

$X$ et $Y$ sont donc les racines de $t^{2}-\dfrac72 t + \dfrac32 = 0$,
c'est-à-dire de $2t^{2}-7t+3 = 0$, de discriminant $25$ : $t = 3$ ou
$t = \dfrac12$.

Les couples solutions sont $(x;y) = \left(e^{3}\,;\sqrt{e}\right)$ et
$\left(\sqrt{e}\,;e^{3}\right)$.

\smallskip
\textbf{L'intégrale.} On écrit $x+3 = (x+1)+2$, d'où
\[
  J = \int_{1}^{3}\frac{\dx}{(x+1)^{2}} + 2\int_{1}^{3}\frac{\dx}{(x+1)^{3}}
    = \left[-\frac{1}{x+1}-\frac{1}{(x+1)^{2}}\right]_{1}^{3},
\]
puis
\[
  J = \left(-\frac14-\frac1{16}\right)-\left(-\frac12-\frac14\right)
    = \frac{7}{16}.
\]
\end{corrige}


% =========================================================================
\begin{devoir}[2 heures]
Trois exercices indépendants, notés sur $20$. Le devoir est cumulatif : il porte sur ce chapitre et sur ceux qui le précèdent. Les énoncés sont repris de sessions réelles et assemblés ici en épreuve ; leur provenance est indiquée à chaque fois.

\exods{1}{D'après le baccalauréat probatoire, Cameroun, série C, session 2021 — exercice 3}{5 points}
Soit $f$ définie sur $\R\setminus\{1\}$ par $f(x) = \dfrac{x^{2}-4}{1-x}$, de
courbe $(C)$.
\begin{enumerate}[label=\textbf{\alph*)}, leftmargin=*, itemsep=0.15em]
  \item Déterminer trois réels $a$, $b$ et $c$ tels que
        $f(x) = ax+b+\dfrac{c}{x-1}$.
  \item En déduire que la droite $(T)$ d'équation $y = -x-1$ est asymptote
        oblique à $(C)$, et étudier la position de $(C)$ par rapport à $(T)$.
  \item Construire $(C)$, son asymptote verticale et $(T)$ dans un même
        repère.
\end{enumerate}

\exods{2}{Concours d'entrée, ENI Fianarantsoa, session 2008-2009 — exercice 3, suite}{7 points}
En déduire la valeur de
$\displaystyle J = \int_{1}^{3}\frac{x+3}{(x+1)^{3}}\dx$.

\medskip

\exods{3}{Baccalauréat, Mauritanie, série C, session 2022 — exercice 4}{8 points}
Sur $\intervalleo{0}{+\infty}$ on pose $f(x) = x-3-\dfrac{3\ln x}{x}$ et
$u(x) = x^{2}-3+3\ln x$.
\begin{enumerate}[label=\textbf{\arabic*)}, leftmargin=*, itemsep=0.15em]
  \item Dresser le tableau de variation de $u$, montrer que $u(x) = 0$ admet
        une solution unique $\alpha$ avec $1{,}4 < \alpha < 1{,}41$, puis
        donner le signe de $u$.
  \item Calculer les limites de $f$ en $0^{+}$ et en $+\infty$, montrer que la
        droite $(D) : y = x-3$ est asymptote et préciser la position relative.
  \item Sachant que $f'(x) = \dfrac{u(x)}{x^{2}}$, dresser le tableau de
        variation de $f$ et construire sa courbe.
  \item Soit $\varphi$ la restriction de $f$ à $\intervalleof{0}{\alpha}$.
        Montrer que $\varphi$ est une bijection de $\intervalleof{0}{\alpha}$
        sur un intervalle $J$ à déterminer, puis construire la courbe de
        $\varphi^{-1}$.
\end{enumerate}

\end{devoir}

\begin{corrige}[devoir surveillé]
\textbf{Exercice 1. a)} On réduit au même dénominateur :
$-x-1+\dfrac{3}{x-1} = \dfrac{-(x+1)(x-1)+3}{x-1}
= \dfrac{-x^{2}+4}{x-1} = \dfrac{x^{2}-4}{1-x}$. Donc $a = -1$, $b = -1$ et
$c = 3$.
\textbf{b)} $f(x)-(-x-1) = \dfrac{3}{x-1}$, qui tend vers $0$ aux deux
infinis : la droite $(T)$ est asymptote oblique. Cette différence est
strictement positive pour $x>1$ et strictement négative pour $x<1$ : la courbe
est au-dessus de $(T)$ à droite de $1$, au-dessous à gauche.
\textbf{c)} La courbe possède deux branches, séparées par l'asymptote
verticale $x = 1$ ; toutes deux décroissent, l'une venant de $+\infty$ et
descendant vers $-\infty$ le long de $x = 1$, l'autre repartant de $+\infty$
au-dessus de $(T)$ pour descendre vers $-\infty$.

\textbf{Exercice 2.} Chacun des deux morceaux est de la forme $u'u^{n}$ avec
$u = x+1$ :
\[
  \int_{1}^{3}\frac{\dx}{(x+1)^{2}} = \left[\frac{-1}{x+1}\right]_{1}^{3}
  = -\frac14+\frac12 = \frac14,
  \qquad
  \int_{1}^{3}\frac{2\dx}{(x+1)^{3}} = \left[\frac{-1}{(x+1)^{2}}\right]_{1}^{3}
  = -\frac{1}{16}+\frac14 = \frac{3}{16}.
\]
En additionnant, $J = \dfrac14+\dfrac{3}{16} = \dfrac{7}{16}$.

\textbf{Exercice 3.} $u'(x) = 2x+\frac3x>0$ : $u$ croît de $-\infty$ à $+\infty$, donc
s'annule une fois et une seule. Comme $u(1{,}4)<0<u(1{,}41)$, la racine
$\alpha$ est dans $\intervalleo{1{,}4}{1{,}41}$ ; $u$ est négative avant
$\alpha$, positive après. En $0^{+}$, $\frac{\ln x}{x}\to-\infty$ donc
$f(x)\to+\infty$ ; en $+\infty$, $\frac{\ln x}{x}\to0$ et
$f(x)-(x-3) = -\frac{3\ln x}{x}\to0$ : la droite $y = x-3$ est asymptote, la
courbe est au-dessus pour $x<1$ et en dessous pour $x>1$. Le signe de $f'$
étant celui de $u$, $f$ décroît sur $\intervalleof{0}{\alpha}$ et croît
ensuite. Enfin $\varphi$ est continue et strictement décroissante de
$\intervalleof{0}{\alpha}$ sur $J = \intervallefo{f(\alpha)}{+\infty}$ :
c'est une bijection, et la courbe de $\varphi^{-1}$ s'obtient par symétrie
par rapport à la droite $y = x$.

\end{corrige}
